% SPDX-License-Identifier: Apache-2.0
% covers: Hdmi
\datasheet{Hdmi}{The video peripheral for an HDMI encoder chip: a raster, a framebuffer, and AXI-Lite to paint it.}

\dsfacts{\code{lib/parts/src/hdmi.rs}}{\code{txhdl\_parts::hdmi::Hdmi}}%
{\code{//docs:hdmi}}

\subsection*{Function}

\code{Hdmi} drives the parallel inputs of an encoder chip such as the
SiI9134 on the Alinx AX7A200B. It is one unit on the pixel clock. It
counts the raster, reads its framebuffer under the beam, and drives
\code{rgb}, 24 bits with red on top, \code{hsync}, \code{vsync} and
\code{de} from registers. A host paints the framebuffer through the
unit's AXI-Lite side: \code{aw}, \code{ar} and \code{w} in, \code{b}
and \code{r} out, with 32-bit addresses and words.

\subsection*{Parameters}

\begin{center}\footnotesize
\begin{tabular}{@{}lp{0.7\textwidth}@{}}
\toprule
Parameter & Meaning \\
\midrule
\code{HV} & Visible columns. \\
\code{HFP} & Columns of horizontal front porch. \\
\code{HSW} & Columns of horizontal sync. \\
\code{HBP} & Columns of horizontal back porch. \\
\code{VV} & Visible rows. \\
\code{VFP} & Rows of vertical front porch. \\
\code{VSW} & Rows of vertical sync. \\
\code{VBP} & Rows of vertical back porch. \\
\code{SHIFT} & A framebuffer pixel covers \code{1 << SHIFT} by
\code{1 << SHIFT} pixels of the screen. \\
\bottomrule
\end{tabular}
\end{center}

The tables below are for \code{Hdmi<640, 16, 96, 48, 480, 10, 2, 33,
2>}. Those are the constants of \code{hdmi::vga}, 640 by 480 at 60~Hz
as VESA states it, with \code{SHIFT} of 2. That is the mode of
\code{hdmi/src/netlist.rs}, and its picture is 160 by 120.

\subsection*{Register map}

The unit decodes the word from address bits 3 to 2, so the words are at
offsets \code{0x0}, \code{0x4} and \code{0x8} of its range.
\code{ex\_hdmi} puts them at \code{0x1000}, \code{0x1004} and
\code{0x1008}.

\begin{center}\footnotesize
\begin{tabular}{@{}lllp{0.5\textwidth}@{}}
\toprule
Word & Offset & Access & Bits \\
\midrule
0, status & \code{0x0} & read & 0: high in vertical blanking. 31 to 16:
the count of frames shown. \\
1, cursor & \code{0x4} & read, write & 7 to 0: the column. 14 to 8: the
row. \\
2, pixel & \code{0x8} & write & 11 to 0: the colour, four bits each of
red, green and blue from the top. \\
\bottomrule
\end{tabular}
\end{center}

\dstables{Hdmi}

\subsection*{Behaviour}

\begin{itemize}
\item \code{hc} counts the column and \code{vc} the row, each over the
visible part and the blanking. \code{frames} goes up by one at the end
of each frame.
\item At each edge the unit reads the framebuffer at the pixel under
the beam into \code{px}. \code{//docs:hdmi} calls this the synchronous
read a block RAM has. The framebuffer address is the row,
\code{vc >> SHIFT}, in the top seven bits and the column,
\code{hc >> SHIFT}, in the low eight.
\item At the same edge the unit registers the syncs in \code{hs\_q} and
\code{vs\_q}, low inside the sync pulse, and the enable in
\code{de\_q}, high over the visible part. So they meet the pixel they
belong to.
\item \code{rgb} is \code{px} with each 4-bit colour widened to 8 bits
by repeating its bits, so 15 is full scale.
\item A write to the cursor word sets \code{cx} and \code{cy}.
\item A write to the pixel word puts the colour into the framebuffer at
the cursor, and moves the cursor to the next column. From the last
column it moves to the first column of the next row, and from the last
row to the first. A fixed burst to this word paints a run of pixels.
\item A write goes when \code{aw} and \code{w} both offer and \code{b}
is ready. A read goes when \code{ar} offers and \code{r} is ready. The
unit answers every transaction \code{Okay}. A read of word 2 or word 3
answers zero, and a write to word 0 or word 3 does nothing else.
\end{itemize}

\subsection*{Verification}

\begin{itemize}
\item \code{the\_raster\_is\_where\_the\_mode\_puts\_it}, in
\code{//lib/parts:parts\_test}, runs \code{Hdmi<8, 2, 3, 3, 6, 1, 2, 2,
1>} for two frames. It checks that exactly one lag of the outputs
behind the counters fits \code{de}, \code{hsync} and \code{vsync}
together.
\item \code{ex\_hdmi} runs \code{Hdmi<16, 2, 3, 3, 12, 1, 2, 2, 1>}
behind \code{LiteBridge1} and an \code{AxiHost}. The host sets the
cursor to the corner and paints all 48 pixels of the 8 by 6
framebuffer in one fixed burst. The run rebuilds the last whole frame
from the pixels \code{de} marks, and asserts every one of its 192
screen pixels against the colour painted.
\item The build simulates that unit's netlist against the run under
nvc and Verilator: \code{//docs:hdmi\_sim\_hdmi\_video\_hdmi\_video\_tb\_test}
and \code{//docs:hdmi\_vsim\_hdmi\_video\_test}. The co-simulated mode
is the example's, not the one in the tables above.
\end{itemize}

\subsection*{Used by}

The example \code{ex\_hdmi}. The demonstration for the AX7A200B:
\code{//hdmi:netlist} writes the Verilog of \code{hdmi\_video} in the
640 by 480 mode, with a test picture in the framebuffer's initial
values. \code{hdmi/board/hdmi\_demo.v} instantiates it on a 25.2~MHz
pixel clock, with its AXI-Lite inputs held idle.
\code{//hdmi:demo\_synth} and \code{//hdmi:demo\_pnr} put the design
through Vivado to a bitstream. \code{//docs:hdmi} states that it has
not yet run on the board.

\subsection*{Limits}

A framebuffer pixel is 12 bits, and the framebuffer holds
\code{FB\_WORDS}, $2^{15}$ words, so the picture is at most 256
columns by 128 rows. \code{hc} and \code{vc} are 12 bits wide. The mode
is fixed when the type is named. The whole unit runs on the pixel
clock, bus included. A host on another clock needs a crossing of the
five AXI-Lite channels, and nothing in the tree provides one. The
framebuffer cannot be read over the bus. On the board, \code{//docs:hdmi}
marks two things as unconfirmed: which of the chip's data inputs are
red, green and blue, and the I/O standard of the chip's pins.
